\documentclass[twocolumn,12pt]{article}
\usepackage{lipsum}
\usepackage{graphicx}
\usepackage{amsmath}
\usepackage[hidelinks]{hyperref}

\title{The Epistemology of Computational Systems}
\author{A.~R.~Fontaine \and M.~J.~Hessler}
\date{2024}

\begin{document}

\maketitle

\begin{abstract}
This paper examines the epistemological foundations of modern computational
systems, focusing on the boundary conditions under which formal verification
methods yield reliable knowledge claims.
We argue that Gettier-style counterexamples apply with equal force to
machine-generated proof certificates.\footnote{The term ``proof certificate''
is used here in the sense of Necula \& Lee (1998), not the informal
usage common in software engineering.}
A taxonomy of computational epistemic failures is proposed, grounded in
an analysis of three representative system classes.
\end{abstract}

\section{Introduction}
\label{sec:intro}

Computational epistemology sits at the intersection of formal methods,
philosophy of science, and cognitive systems theory.
\lipsum[1]
A defining challenge\footnote{This challenge was first articulated in the
context of automated theorem provers by Robinson (1965), though he did not
use this terminology.} is distinguishing between \emph{verified} results
and \emph{justified} results.

\lipsum[2]

The distinction matters in practice: a system may produce a verified proof
of a false theorem if the formal model diverges from the intended
semantics.\footnote{The canonical example is the Pentium FDIV bug (1994),
in which a formally verified microcode routine produced incorrect results
for a specific class of floating-point inputs.} We return to this point
in Section~\ref{sec:taxonomy}.

\subsection{Scope and Method}
\label{sec:scope}

The analysis is restricted to deterministic, finitely-computable systems.
Probabilistic and quantum computational models are deferred to future work.
\lipsum[3]

\subsubsection{Formal Preliminaries}
\label{sec:prelim}

Let $\Sigma$ be a finite alphabet and $L \subseteq \Sigma^*$ a decidable language.
A \emph{verifier} $V$ is a deterministic Turing machine such that for all
$x \in L$, there exists a certificate $c$ with $V(x,c) = 1$.
\lipsum[4]

\section{A Taxonomy of Epistemic Failure Modes}
\label{sec:taxonomy}

We identify four primary failure modes in computational knowledge systems.
Each mode corresponds to a distinct breakdown in the justification chain
from axioms to conclusions.

\subsection{Type I: Model-Semantic Divergence}
\label{sec:type1}

Model-semantic divergence occurs when the formal model $M$ is a sound
abstraction of the specification $S$ but $S$ is itself an inaccurate
description of the intended system $\mathcal{I}$.
\lipsum[5]

\subsubsection{Illustrative Cases}

The TLS 1.3 handshake verification effort \cite{bhargavan2017} provides
a well-documented instance of this failure mode.\footnote{Bhargavan et
al.\ verified the cryptographic core of TLS 1.3 using F$\star$, but noted
that their model excluded session resumption tickets—an exclusion that
proved non-trivial in later implementations.}
\lipsum[6]

\subsection{Type II: Proof-State Corruption}
\label{sec:type2}

Proof-state corruption arises when a proof assistant's internal state is
modified by a trusted kernel extension in a way that is not reflected in
the proof certificate.
\lipsum[7]

\begin{figure}[h]
  \centering
  \rule{0.9\columnwidth}{3cm}
  \caption{Schematic of the justification chain from axioms ($\mathcal{A}$)
  through the proof kernel ($K$) to the certificate ($C$). A Type~II failure
  corrupts the $K \to C$ link without invalidating $\mathcal{A} \to K$.}
  \label{fig:justification-chain}
\end{figure}

\lipsum[8]

\subsection{Type III: Certificate Non-Portability}
\label{sec:type3}

A certificate $c$ produced by verifier $V_1$ may not be accepted by verifier
$V_2$ even when both are sound with respect to the same logic, because the
two systems may use non-equivalent encodings of the proof term.
\lipsum[9]

\subsubsection{Cross-System Verification}

Cross-system verification remains an open problem.\footnote{The QED Manifesto
(1994) proposed a universal proof repository, but the project stalled
precisely because of certificate non-portability across the Coq, Isabelle,
and HOL family of proof assistants.}
\lipsum[10]

\section{Implications for Verified Software Engineering}
\label{sec:implications}

The taxonomy developed in Section~\ref{sec:taxonomy} has direct implications
for the deployment of formally verified components in safety-critical systems.

\subsection{Certification Standards}

Current certification standards such as DO-178C (avionics) and IEC 61508
(industrial safety) do not yet account for Type~II or Type~III failures.
\lipsum[11]

\subsubsection{Proposed Augmentations}

We propose three augmentations to existing standards.
First, certificate portability testing should be mandated for all Level~A
software.\footnote{Level A software is defined by DO-178C as software whose
failure would result in catastrophic failure conditions.}
Second, formal models should include explicit semantic gap documentation.
Third, trusted kernel extensions should undergo independent formal verification.

\lipsum[12]

\subsection{Tool Qualification}

Tool qualification under DO-330 partially addresses Type~I failures but
does not cover the certificate layer.
\lipsum[13]

\section{Conclusion}
\label{sec:conclusion}

We have presented a taxonomy of four epistemic failure modes in computational
verification systems and argued that existing certification standards are
insufficient to address the full range of failure modes.
\lipsum[14]

\begin{thebibliography}{9}

\bibitem{bhargavan2017}
K.~Bhargavan, B.~Blanchet, and N.~Kobeissi.
\newblock Verified models and reference implementations for the {TLS} 1.3
  standard candidate.
\newblock In \emph{Proceedings of the IEEE Symposium on Security and Privacy},
  pages 483--502, 2017.

\bibitem{necula1998}
G.~C.~Necula and P.~Lee.
\newblock Safe kernel extensions without run-time checking.
\newblock In \emph{Proceedings of the 2nd USENIX Symposium on Operating Systems
  Design and Implementation}, 1998.

\bibitem{robinson1965}
J.~A.~Robinson.
\newblock A machine-oriented logic based on the resolution principle.
\newblock \emph{Journal of the ACM}, 12(1):23--41, 1965.

\end{thebibliography}

\end{document}
